% Faithful official CCC2022 article16 additional pKt coding/approximation and agnostic-learning definitions.
\section*{2.6 Agnostic learning model (author definitions)}
Consider the problem of learning an unknown concept $f:X\to\{0,1\}$ over a finite domain $X$. We will always let $X=\{0,1\}^n$, for some $n\geq1$. A concept class $\mathcal C$ is a collection of concepts of this form. For a fixed $n$, we let $\mathcal C_n$ denote $\mathcal C\cap\{f:\{0,1\}^n\to\{0,1\}\}$. For a size function $s:\mathbb N\to\mathbb N$, we let
\[\mathrm{SIZE}[s]=\{f:\{0,1\}^n\to\{0,1\}\mid n\in\mathbb N\text{ and }f\text{ admits a size-}s(n)\text{ circuit}\}\]
denote the concept class of Boolean circuits of size (number of gates) at most $s(n)$ over $n$ input variables. For definiteness, we consider circuits over AND and OR gates of fan-in $2$ and NOT gates.
A randomized Boolean function $f$ maps an input $x$ to a distribution $D$ supported over $\{0,1\}$. Given a distribution $D_n$ supported over $\{0,1\}^n$, a function $h:\{0,1\}^n\to\{0,1\}$, and a (possibly randomized) Boolean function $f:\{0,1\}^n\to\{0,1\}$, we let
\[\operatorname{err}_{D_n}(h,f)=\Pr_{f,x\sim D_n}[h(x)\neq f(x)].\]
We will consider the (agnostic) learnability of a concept class $\mathcal C$ with respect to a class $\mathcal D$ of ensembles $D=\{D_n\}_{n\geq1}$ of distributions $D_n$. We use $\mathcal D_n$ to denote $\{D_n\mid D\in\mathcal D\}$. For simplicity, we refer to $\mathcal D$ as a class of distributions.
Given a randomized function $f:\{0,1\}^n\to\{0,1\}$ and a distribution $D_n$, the learning algorithm has access to an example oracle $\mathrm{EX}(f,D_n)$ that behaves as follows: each query to $\mathrm{EX}(f,D_n)$ returns an independent and identically distributed pair $(x,b)$, where $x$ is sampled according to $D_n$ and $b=f(x)$. The sample complexity of the learning algorithm is the number of queries made to $\mathrm{EX}(f,D_n)$.
Before formalising the agnostic learning model, we provide an informal overview. The learning algorithm knows the class $\mathcal D$ of distributions and the concept class $\mathcal C$. It is given access to $\mathrm{EX}(f,D_n)$ for some unknown $D\in\mathcal D_n$ and an arbitrary (possibly randomized) function $f:\{0,1\}^n\to\{0,1\}$. The goal of the learning algorithm is to produce, with high probability over its internal randomness and queries to $\mathrm{EX}(f,D)$, a hypothesis $h$ such that $\operatorname{err}_D(h,f)\leq\operatorname{opt}_{\mathcal C_n,D,f}+\varepsilon$, where
\[\operatorname{opt}_{\mathcal C_n,D,f}=\min_{c\in\mathcal C_n}\operatorname{err}_D(c,f).\]
In other words, the learning algorithm is required to output an efficient representation of a hypothesis $h:\{0,1\}^n\to\{0,1\}$ that is almost as accurate as the best concept in $\mathcal C_n$ with respect to the pair $(D,f)$. (Note that $h$ is not required to be a function in $\mathcal C$.)
\paragraph{Definition 15 (Agnostic PAC learning [20]).} Let $\mathcal C$ be a concept class, and let $\mathcal D$ be a class of distributions. We say that a randomized algorithm $A$ agnostically learns $\mathcal C$ on $\mathcal D$ if the following holds. For every $n\geq1$, distribution $D\in\mathcal D_n$, and randomized function $f:\{0,1\}^n\to\{0,1\}$, when $A$ is given oracle access to $\mathrm{EX}(f,D)$ and receives as input $n,\varepsilon>0$, and $\delta>0$,
\[\Pr_{A,\mathrm{EX}(f,D)}[A^{\mathrm{EX}(f,D)}\text{ outputs a hypothesis }h\text{ such that }\operatorname{err}_D(h,f)\leq\operatorname{opt}_{\mathcal C_n,D,f}+\varepsilon]\geq1-\delta.\]
We remark that an equivalent way of formulating the agnostic learning model is by considering distributions $D$ supported over $\{0,1\}^{n+1}$. In this case, there is no need to refer to randomized functions, and one measures the error of a hypothesis $h$ via $\operatorname{err}_D(h)=\Pr_{(x,b)\sim D}[h(x)\neq b]$, where $x\in\{0,1\}^n$ and $b\in\{0,1\}$.
\section*{Additional author pKt properties}
\paragraph{Lemma 18.} There is a universal constant $c>0$ such that the following holds. For every time bound $t\in\mathbb N$ and $x\in\{0,1\}^n$,
\[\mathrm K(x\mid t)\leq\mathrm{pK}^t(x)+c\log n.\]
\paragraph{Proof.} Let $\ell:=\mathrm{pK}^t(x)$. Since $x$ is of length $n$, we can represent $\ell$ using $O(\log n)$ bits. Given $t,\ell$ and $n$, it is possible to sample $x$ with probability $\mu\geq\frac23\cdot2^{-\ell}$ by randomly guessing $w\sim\{0,1\}^t$ and $M\sim\{0,1\}^\ell$ then simulating $M(w)$ for $t$ steps. Consequently, by the coding theorem for (time-unbounded) Kolmogorov complexity [24], we get that $\mathrm K(x\mid t)\leq\log(1/\mu)+O(\log n)\leq\mathrm{pK}^t(x)+O(\log n)$.\hfill$\triangleleft$
\paragraph{Lemma 20 (Probabilistic Incompressibility).} For any string $y\in\{0,1\}^*$, time bound $t\in\mathbb N$ (including $t=\infty$), $\delta\in(0,1]$, and positive integer $\alpha$, we have
\[\Pr_{x\sim\{0,1\}^n}(\mathrm{pK}_\delta^t(x\mid y)\leq n-\alpha)\leq\frac1{2^{\alpha-1}\cdot\delta}.\]
\paragraph{Proof.} For the sake of contradiction, suppose the statement of the lemma does not hold. Then by the definition of $\mathrm{pK}_\delta^t$, we have
\begin{multline*}
\mathbb E_{x\sim\{0,1\}^n,w\sim\{0,1\}^t}[\exists M_{(x,w)}\in\{0,1\}^{\leq n-\alpha}\text{ such that }M_{(x,w)}(w,y)\text{ outputs }x\text{ within }t\text{ steps}]\\
>1/2^{\alpha-1}.
\end{multline*}
By averaging, there is a way to fix $w\in\{0,1\}^t$ that at least preserves this expectation, yielding
\begin{multline*}
\Pr_{x\sim\{0,1\}^n}[\exists M_{(x,w)}\in\{0,1\}^{\leq n-\alpha}\text{ such that }M_{(x,w)}(w,y)\text{ outputs }x\text{ within }t\text{ steps}]\\
>1/2^{\alpha-1}.
\end{multline*}
However, by counting, the above probability is at most $2^{n-\alpha+1}/2^n=1/2^{\alpha-1}$. A contradiction.\hfill$\triangleleft$
\paragraph{Lemma 24 (Unconditional Source Coding for $\mathrm{pK}^t$ [34]).} There exists a polynomial $p$ such that for any $T,a:\mathbb N\to\mathbb N$, $n\in\mathbb N$, $D_n\in\mathrm{Samp}[T(n)]/a(n)$, and $x\in\operatorname{Supp}(D_n)$,
\[\mathrm{pK}^{p(T(n))}(x)\leq\log(1/D_n(x))+O(\log(T(n)))+a(n).\]
For a distribution $D_n$ supported over $X$ and an integer $m\geq1$, we let $D_n^m:=D_n\times\ldots\times D_n$ be the probability distribution supported over $X^m$ obtained by taking the product of $m$ independent copies of $D_n$.
As a consequence of Lemma 24, we obtain the following result.
\paragraph{Lemma 25 (Source Coding for $m$ copies of $D_n$).} There exists a polynomial $p$ such that for any $T,a:\mathbb N\to\mathbb N$, $n,m\in\mathbb N$, $D_n\in\mathrm{Samp}[T(n)]/a(n)$, and $x\in\operatorname{Supp}(D_n^m)$,
\[\mathrm{pK}^{p(T(n),m)}(x)\leq\log(1/D_n^m(x))+O(\log(T(n)))+a(n)+O(\log(m)).\]
\paragraph{Proof.} Since $D_n$ can be sampled in time $T(n)$ using $a(n)$ bits of advice, it is possible to sample from $D_n^m$ in time $\operatorname{poly}(T(n),m)$ using $a(n)+O(\log m)$ bits of advice, where the extra advice is used to encode $m$. The result now immediately follows from Lemma 24.\hfill$\triangleleft$
\subsection*{3.5 Approximating $\mathrm{pK}^t$ under randomized average-case easiness}
\paragraph{Lemma 27.} If $\mathrm{DistNP}\subseteq\mathrm{AvgBPP}$, then there exists a polynomial $\tau$ such that the following promise problem is in promiseBPP:
\begin{align*}
\Pi_{\mathrm{YES}}&:=\{(x,1^s,1^t)\mid\mathrm{pK}^t(x)\leq s\},\\
\Pi_{\mathrm{NO}}&:=\{(x,1^s,1^t)\mid\mathrm{pK}^{\tau(t)}(x)>s+\log\tau(t)\}.
\end{align*}
\paragraph{Proof.} Let $k=s+\log n$. Consider the language
\[L:=\{(\mathrm{DP}_k(x;z),w,1^s,1^n)\mid\exists M\in\{0,1\}^s,M(w)\text{ outputs }x\in\{0,1\}^n\text{ within }|w|\text{ steps}\}.\]
Note that $L\in\mathrm{NP}$. Define a distribution family $\mathcal D=\{D_{\langle nk+k,t,s,n\rangle}\}$, each member of which does the following: sample $u\sim U_{nk+k}$ and $w\sim U_t$, and then output $(u,w,1^s,1^n)$. By assumption, $(L,\mathcal D)\in\mathrm{AvgBPP}$. Let $B$ be a randomized heuristic algorithm for $(L,\mathcal D)$ as described in Lemma 6.
Now, define an algorithm $B'$:
\begin{itemize}
\item On input $(x,1^s,1^t)$ with $x\in\{0,1\}^n$, sample $z\sim U_{nk}$ and $w\sim U_t$, and then output $B(\mathrm{DP}_k(x;z),w,1^s,1^n)$.
\end{itemize}
Below, we show that $B'$ solves $(\Pi_{\mathrm{YES}},\Pi_{\mathrm{NO}})$ correctly with high probability in the worst case.
First, consider the case that $(x,1^s,1^t)\in\Pi_{\mathrm{YES}}$. By the definitions of $L$ and pK, for any choice of $z\in\{0,1\}^{nk}$,
\[\Pr_w[(\mathrm{DP}_k(x;z),w,1^s,1^n)\in L]\geq2/3.\]
The definition of $B$ then implies that
\[\Pr_{w,z,r_B}[B(\mathrm{DP}_k(x;z),w,1^s,1^n)=1]>1/2,\]
and so
\[\Pr_{r_{B'}}(B'(x,1^s,1^t)=1)>1/2,\tag{13}\]
where $r_B$ denotes the internal randomness of $B$, and $r_{B'}=(w,z,r_B)$ that of $B'$.
Now consider the case that $(x,1^s,1^t)\in\Pi_{\mathrm{NO}}$. For a contradiction, suppose that
\[\Pr_{w,z,r_B}[B(\mathrm{DP}_k(x;z),w,1^s,1^n)=1]>1/3.\tag{14}\]
By a counting argument, for randomly selected $u$ and $w$,
\[\Pr_{u,w}[(u,w,1^s,1^n)\in L]\leq\frac{2^s\cdot2^{nk}\cdot2^{|w|}}{2^{nk+k+|w|}}=1/n,\]
where the last line follows from the definition of $k=s+\log n$. Then by definition of $B$,
\[\Pr_{u,w,r_B}[B(u,w,1^s,1^n)=1]\leq4/n=o(1).\tag{15}\]
Comparing Equations (14) and (15), it is clear that $B(-,U_t,1^s,1^n)$ $(1/4)$-distinguishes $\mathrm{DP}_k(x;U_{nk})$ from $U_{nk+k}$. Lemma 22 implies that
\[\mathrm{pK}^{p'(t)}(x)\leq k+\log p'(t)=s+\log n+\log p'(t),\]
for some polynomial $p'$ with $p'(t)\geq p_{\mathrm{DP}}(t_B\cdot4\cdot n)$, where $p_{\mathrm{DP}}$ is the polynomial from Lemma 22 and $t_B$ denotes the time required to compute $B(-,U_t,1^s,1^n)$. For the polynomial $\tau(t)=t\cdot p'(t)$, this means that $(x,1^s,1^t)$ is not in $\Pi_{\mathrm{NO}}$, which gives the desired contradiction. By definition of $B'$, we have that
\[\Pr_{r_{B'}}(B'(x,1^s,1^t)=1)\leq1/3.\tag{16}\]
By Equations (13) and (16), $B'$ yields a promiseBPP algorithm for $(\Pi_{\mathrm{YES}},\Pi_{\mathrm{NO}})$ via standard success amplification techniques.\hfill$\triangleleft$
\paragraph{Lemma 28 (Approximating $\mathrm{pK}^t$).} If $\mathrm{DistNP}\subseteq\mathrm{AvgBPP}$, then there exist a polynomial $\tau$, a constant $C\geq1$, and a randomized algorithm $\mathrm{Approx}_{\tau}\text{-}\mathrm{pK}$ that, on input $(x,1^t)$ where $x\in\{0,1\}^n$ and $t\geq Cn$, runs in time $\operatorname{poly}(n,t)$ and with probability at least $1-o(1)$ outputs an integer $\widetilde s$ such that
\[\mathrm{pK}^{\tau(t)}(x)-\log\tau(t)\leq\widetilde s\leq\mathrm{pK}^t(x).\]
\paragraph{Proof.} Consider a polynomial-time randomized algorithm $A$ that solves the promise problem from Lemma 27. Assume without loss of generality that on the inputs satisfying the promise its error is at most $1/n^2$, where $n=|x|$. Algorithm $\mathrm{Approx}_\tau\text{-}\mathrm{pK}$ runs $A$ on $(x,1^s,1^t)$ for $s=1,2,\ldots,n+\log n$, and outputs the first $\widetilde s$ such that $A(x,1^{\widetilde s},1^t)=1$.
The correctness of $\mathrm{Approx}_\tau\text{-}\mathrm{pK}$ follows by a union bound. Indeed, if $s<\mathrm{pK}^{\tau(t)}(x)-\log\tau(t)$, i.e., $\mathrm{pK}^{\tau(t)}(x)>s+\log\tau(t)$, using the promise we get that $\Pr_A[A(x,1^s,1^t)=1]\leq1/n^2$. On the other hand, if $s=\mathrm{pK}^t(x)$, which implies that $\mathrm{pK}^t(x)\leq s$ and the promise is satisfied, we have $\Pr_A[A(x,1^s,1^t)=1)]\geq1-1/n^2$. Since $\mathrm{pK}^t(x)\leq n+\log n$ if $t\geq Cn$, where $C$ is a sufficiently large constant, with high probability over the internal randomness of $\mathrm{Approx}_\tau\text{-}\mathrm{pK}$, it outputs a value $\widetilde s$ such that $\mathrm{pK}^{\tau(t)}(x)-\log\tau(t)\leq\widetilde s\leq\mathrm{pK}^t(x)$.\hfill$\triangleleft$
